\question[3]
Die Liste a enthält die Zahlen 3 7 5. Die Schleife soll alle Elemente ausgeben.
\begin{lstlisting}
a.reset()
while not a.endpos():
    a.advance()
    print(a.elem())
\end{lstlisting}
\begin{parts}
\part Was geschieht stattdessen?
\part Korrigiere die Schleife.
\end{parts}
\begin{solutionbox}{4.25cm}
a) Ausgegeben werden 7 und 5, die 3 wird übersprungen. Nach der 5 ist das Ende erreicht und \texttt{a.elem()} wirft
einen RuntimeError. \\
b) Die beiden Zeilen im Schleifenrumpf werden getauscht:
\begin{lstlisting}
a.reset()
while not a.endpos():
    print(a.elem())
    a.advance()
\end{lstlisting}
\end{solutionbox}
